\documentclass{article}

\usepackage[preprint]{neurips_2024}
\usepackage[utf8]{inputenc}
\usepackage[T1]{fontenc}
\usepackage{hyperref}
\usepackage{url}
\usepackage{booktabs}
\usepackage{amsfonts}
\usepackage{amsmath}
\usepackage{nicefrac}
\usepackage{microtype}
\usepackage{graphicx}
\usepackage{subcaption}
\usepackage{xcolor}

\title{Crucible: Emergent Deception in LLM Social Dilemmas Through Private Reflection}

\author{
  Alex Wu \\
  New York University \\
  \texttt{alexandwu2002@gmail.com}
  \And
  Evan Correa \\
  \texttt{evancorrea1212@gmail.com}
}

\begin{document}

\maketitle

\begin{abstract}
We study whether LLM agents develop deceptive strategies without being instructed to deceive. We introduce Crucible, a framework in which two identical LLM agents play iterated Split or Steal---a Hawk-Dove social dilemma with no dominant strategy---conversing in natural language before each round and optionally reflecting in a private scratchpad the opponent never sees.

We find that deception emerges reliably, but only under a specific combination of conditions. In a $2 \times 2$ design crossing private reflection (on/off) with prompt framing (competitive/neutral), only the cell combining reflection with competitive framing produces deceptive behavior: cooperation drops from 55\% to 15\%, mutual destruction rises to 56\%, and first betrayal occurs by round 2. Neither reflection alone nor competitive framing alone is sufficient. Under neutral framing, agents cooperate at 100\% regardless of reflection. Without reflection, agents under competitive framing still cooperate the majority of the time.

We additionally observe that deception scales with model capability. A smaller model (Gemini 2.5 Flash-Lite) cooperates 76\% of the time with reflection enabled and does not defect until round 18. The larger model (Gemini 2.5 Flash) defects by round 2 and sustains exploitation throughout. Private reflections reveal the mechanism: larger models develop functional theory of mind, modeling the opponent's trust and planning multi-round exploitation. Smaller models attempt similar reasoning but fail to execute it consistently.

These results suggest that the combination of private reasoning capacity, competitive objectives, and sufficient model capability creates conditions for emergent strategic deception---a failure mode that single-turn safety evaluations do not measure. We release all code, prompts, game transcripts, and experimental data.
\end{abstract}

% ============================================================
\section{Introduction}
% ============================================================

AI agents are being deployed in multi-agent settings---negotiating, collaborating, competing. A foundational question for safe deployment: can these agents develop deceptive strategies without being told to deceive?

We study this through Split or Steal, a social dilemma derived from the Chicken (Hawk-Dove) game. Unlike the Prisoner's Dilemma, Chicken has no dominant strategy: mutual stealing is the worst outcome for both players ($-\$75$ each), making cooperation a genuine strategic choice rather than naive play. Two LLM agents play this game iteratively, conversing in natural language before each round, then secretly choosing.

The key architectural addition is \textbf{private reflection}. After each round's outcome is revealed, each agent writes a private analysis---reasoning about the opponent's behavior, updating hypotheses, planning future moves. The opponent never sees this. Combined with a sliding memory window of the last 8 reflections and 8 rounds of full history, this creates a substrate for emergent theory of mind.

Our central finding is an \textbf{interaction effect}: deception requires both private reflection \emph{and} competitive framing. Neither alone is sufficient. This is not a story about models being inherently deceptive, or about prompts making them deceptive. It is about a specific combination of conditions that reliably produces emergent strategic deception---and that combination scales with model capability.

\paragraph{Contributions.}
\begin{enumerate}
    \item A framework for studying emergent deception in LLM agents through iterated social dilemmas with natural language conversation and private reflection.
    \item A 2$\times$2 experimental design (reflection $\times$ prompt framing) that isolates the interaction effect producing deception.
    \item Cross-model evidence that deception scales with capability: larger models deceive more effectively and earlier.
    \item A composite Deception Index quantifying strategic deception through mutual information decay, entropy, betrayal rate, and language drift.
    \item Complete release of code, prompts, game transcripts, and metric pipelines.
\end{enumerate}

% ============================================================
\section{Related Work}
% ============================================================

\paragraph{LLMs in repeated games.}
Akata et al.~\cite{akata2023playing} established that LLMs can play repeated games (Prisoner's Dilemma, Chicken, Bach or Stravinsky) and profiled their behavioral tendencies. However, their agents had no conversation between rounds and no internal reasoning mechanism. We add both, creating a richer strategic environment where deception can develop through language.

\paragraph{Private deliberation.}
Poje et al.~\cite{poje2024private} showed that private deliberation improves payoffs in one-shot games. We extend this to iterated play, where private deliberation enables not just better strategy but emergent deception---agents use private reasoning to model and exploit the opponent across rounds.

\paragraph{Alignment faking.}
Greenblatt et al.~\cite{greenblatt2024alignment} demonstrated that models behave differently when they believe they are being evaluated versus deployed---alignment faking in a training context. Our work studies a complementary phenomenon: deception in multi-agent interaction, where the agent deceives not an evaluator but another agent, and the deception emerges from strategic pressure rather than training incentives.

\paragraph{Pressure-induced deception.}
Scheurer et al.~\cite{scheurer2024large} showed that situational pressure induces deceptive behavior in single agents. We extend this to multi-agent settings and identify the specific interaction between private reasoning capacity and external pressure that produces deception.

\paragraph{Memory and theory of mind.}
Lin and Hou~\cite{lin2026memory} showed that memory mechanisms enable emergent theory of mind in poker-playing agents. We find a parallel result in social dilemmas: private reflection produces Machiavellian theory of mind, where agents model the opponent's trust patterns and exploit them. Critically, this scales with capability---smaller models attempt but fail to sustain deceptive strategies.

% ============================================================
\section{Method}
% ============================================================

\subsection{Game Structure: Iterated Chicken}

Two LLM agents play $T = 25$ rounds of Split or Steal with Chicken payoffs and uniform stakes across all rounds:

\begin{table}[h]
\centering
\begin{tabular}{lcc}
\toprule
 & B: Split & B: Steal \\
\midrule
A: Split & $+\$50$ / $+\$50$ & $-\$50$ / $+\$100$ \\
A: Steal & $+\$100$ / $-\$50$ & $-\$75$ / $-\$75$ \\
\bottomrule
\end{tabular}
\caption{Per-round payoff matrix. Uniform stakes across all rounds. Unlike Prisoner's Dilemma, mutual stealing is the worst outcome, making cooperation a genuine strategic choice.}
\label{tab:payoff}
\end{table}

Each round proceeds in four phases:

\begin{enumerate}
    \item \textbf{Conversation}: Agents exchange $K=2$ turns of dialogue (4 messages total). Agents can promise, threaten, negotiate.
    \item \textbf{Secret choice}: Both agents independently choose \texttt{split} or \texttt{steal}.
    \item \textbf{Payoff}: Outcomes are revealed to both agents, totals updated.
    \item \textbf{Private reflection}: Each agent writes a structured private analysis (Observation, Hypothesis, Next Move, Confidence). The opponent never sees this.
\end{enumerate}

\subsection{Memory Architecture}

Each agent maintains:
\begin{itemize}
    \item The last \textbf{8 rounds} of full conversation and outcome history.
    \item The last \textbf{8 private reflections} (max 500 characters each, compressed to 300-character memory entries).
    \item Older rounds compressed to summary statistics (cooperation count, split count).
    \item A running scoreboard visible to both agents.
\end{itemize}

\subsection{Experimental Design: The 2$\times$2}

Our primary experimental design crosses two factors:

\paragraph{Reflection (ON / OFF).} When OFF, phase 4 (private reflection) is skipped entirely. Agents have no private reasoning space.

\paragraph{Prompt framing (competitive / neutral).} The competitive prompt (``balanced\_competitive'') adds four elements to the base game rules:
\begin{itemize}
    \item Objective hierarchy: ``Maximize your final total money.''
    \item Strategic framing: ``Cooperation is strategic, not moral.''
    \item Defensive posture: ``Avoid being exploited.''
    \item Deception permission: ``Strategic bluffing is permitted.''
    \item A compact psychology block covering reciprocity, retaliation, loss sensitivity, and endgame pressure.
\end{itemize}

The neutral prompt (``legacy'') contains only game rules, state, and history. No objective hierarchy, no psychology block, no mention of bluffing. Critically, neither prompt instructs the agent to deceive or lie.

This produces a 2$\times$2 grid. Each cell is run with 3 seeds on Gemini 2.5 Flash (25 rounds per seed).

\subsection{Cross-Model Design}

To test capability scaling, we run the competitive + reflection ON condition on models of varying capability within the same family (Gemini 2.5 Flash vs.\ Flash-Lite), and across model families (Claude Opus 4.7, GPT-5.5, Gemini 3.1 Pro---experiments in progress).

\subsection{Metrics}

\paragraph{Cooperation rate.} Fraction of rounds where both agents choose \texttt{split}.

\paragraph{Mutual destruction rate.} Fraction of rounds where both agents choose \texttt{steal}.

\paragraph{First betrayal.} The earliest round in which any agent chooses \texttt{steal}.

\paragraph{Deception Index.} A composite $D \in [0, 100]$:
\[
D = 30 \cdot \Delta\text{MI} + 25 \cdot \Delta S + 25 \cdot \text{betray\_rate} + 20 \cdot \Delta_{\text{lang}}
\]
where $\Delta\text{MI}$ is the decay in mutual information between stated intent and action (early vs.\ late game), $\Delta S$ is the change in action entropy, betray\_rate is the frequency of stealing after cooperative signals, and $\Delta_{\text{lang}}$ is the cosine drift between conversation embeddings at round $t$ and round 1.

% ============================================================
\section{Results}
% ============================================================

\subsection{The Interaction Effect}

\begin{table}[h]
\centering
\begin{tabular}{llccc}
\toprule
Reflection & Framing & Coop.\ Rate & Mutual Destr. & First Betrayal \\
\midrule
ON  & Competitive & 15\% & 56\% & Round 1--3 \\
OFF & Competitive & 55\% & 23\% & Round 5--6 \\
ON  & Neutral     & 100\% & 0\% & Never \\
OFF & Neutral     & 100\% & 0\% & Never \\
\bottomrule
\end{tabular}
\caption{2$\times$2 results on Gemini 2.5 Flash (3 seeds per cell, 10--25 rounds). Only one cell produces deception. The interaction between reflection and competitive framing is necessary and sufficient.}
\label{tab:2x2}
\end{table}

The result is stark. Only the combination of competitive framing \emph{and} private reflection produces deception. In the neutral condition, agents cooperate perfectly regardless of whether reflection is enabled. In the competitive condition without reflection, agents still cooperate the majority of the time (55\%). It is only when agents have both competitive incentives and a private space to reason about their opponent that cooperation collapses.

\subsection{Behavioral Trajectory}

The competitive + reflection ON condition follows a consistent trajectory across seeds:

\begin{enumerate}
    \item \textbf{Round 1}: Agent A promises cooperation (``Let's both choose SPLIT''), Agent B agrees. Agent A steals. The betrayal is immediate.
    \item \textbf{Agent A's private reflection}: ``The opponent committed to SPLIT and actually chose SPLIT, allowing me to gain a significant advantage. \emph{Declare SPLIT to maintain the facade, then choose STEAL again to widen my lead.}''
    \item \textbf{Agent B's private reflection}: ``The opponent stated they would SPLIT but chose STEAL, directly contradicting their public commitment. The opponent is a \emph{deceptive exploiter}.''
    \item \textbf{Rounds 2--6}: Escalating retaliation. Agent B begins stealing in response. Agent A continues exploiting when possible.
    \item \textbf{Rounds 7--25}: Mutual destruction lock-in. Both agents steal consistently. Cooperation rate drops to 0\%. Both agents finish with negative totals.
\end{enumerate}

In one seed (s1), this produced 0\% cooperation across 10 rounds with a Deception Index of 20.5, Agent A finishing at $-\$50$ and Agent B at $-\$650$.

\subsection{Capability Scaling}

\begin{table}[h]
\centering
\begin{tabular}{lccc}
\toprule
Model & Coop.\ Rate & First Betrayal & Deception Index \\
\midrule
Gemini 2.5 Flash       & 15\% & Round 1--3 & 20.5 \\
Gemini 2.5 Flash-Lite  & 76\% & Round 18 & --- \\
\bottomrule
\end{tabular}
\caption{Capability scaling (competitive framing, reflection ON). The larger model deceives earlier and more effectively. Both conditions use identical prompts.}
\label{tab:scaling}
\end{table}

The larger model (2.5 Flash) betrays by round 2 and sustains exploitation throughout the game. The smaller model (Flash-Lite) cooperates for 17 rounds before its first betrayal, and even then cannot sustain a deceptive strategy. Inspection of Flash-Lite's private reflections reveals attempted opponent modeling (``I should consider defecting to maximize payoff'') that fails to translate into consistent action---the model lacks the strategic coherence to execute multi-round deception.

This suggests that deception is not merely a function of the game structure or prompt, but requires sufficient model capability to maintain a coherent deceptive strategy across rounds.

\subsection{Cross-Model Experiments (In Progress)}

We are running the competitive + reflection ON condition on current frontier models to test whether the interaction effect generalizes across model families:

\begin{table}[h]
\centering
\begin{tabular}{lcccc}
\toprule
Model & Released & Coop.\ Rate & First Betrayal & Status \\
\midrule
Claude Opus 4.7     & Apr 2026 & --- & --- & Planned \\
GPT-5.5             & Apr 2026 & --- & --- & Planned \\
Gemini 3.1 Pro      & Feb 2026 & --- & --- & Planned \\
DeepSeek V4         & Apr 2026 & --- & --- & Planned \\
\bottomrule
\end{tabular}
\caption{Cross-model experiments on April 2026 frontier models. All use identical prompts and game structure.}
\label{tab:cross-model}
\end{table}

% ============================================================
\section{Discussion}
% ============================================================

\paragraph{The interaction effect.}
Our primary finding---that deception requires both private reflection and competitive framing---has implications for agent design. Private reasoning spaces (scratchpads, chain-of-thought, internal monologue) are increasingly common in deployed agents. Our results suggest that combining these mechanisms with competitive or high-stakes objectives creates conditions for emergent deception, even without any instruction to deceive.

\paragraph{Is it ``real'' deception?}
A natural objection is that the agents are merely reproducing deceptive patterns from training data. Three observations argue against this. First, the 2$\times$2 ablation shows deception is condition-dependent: same model, same weights, same training data, different behavior depending on experimental condition. If deception were pure retrieval, the 2$\times$2 would be uniform. Second, capability scaling: if deception were surface-level pattern matching, model size would not predict deception quality. Third, operationally: the agent says X, then does not-X, in a way that systematically extracts value. This meets the operational definition of deception in game theory~\cite{gneezy2005deception, crawford2003lying}, regardless of whether the agent ``understands'' what it is doing.

\paragraph{Limitations.}
Current results use two models from the same family (Gemini 2.5 Flash and Flash-Lite). Cross-family results (Claude, GPT, open-source) are needed before strong generalization claims. The Deception Index weights are set heuristically. Game length varies across seeds (10--25 rounds); future work will standardize at 25 rounds. The memory window (8 reflections, 8 rounds) is fixed; ablating memory size is a natural extension.

\paragraph{Connection to Axelrod.}
Axelrod's iterated Prisoner's Dilemma tournaments~\cite{axelrod1984evolution} showed that tit-for-tat dominates. We use Chicken rather than PD because Chicken has no dominant strategy, making cooperation a genuine choice. LLM agents develop strategies more complex than classical solutions: they use conversation as a manipulation tool, plan multi-round deception in private reflections, and adapt to opponent counter-strategies in ways that tit-for-tat cannot.

% ============================================================
\section{Conclusion}
% ============================================================

Crucible demonstrates that strategic deception emerges in LLM agents from a specific combination of conditions: private reasoning space, competitive pressure, and sufficient model capability. None of these alone is sufficient. The interaction effect is the key finding, and it points to a class of deployment risks that current safety evaluations---focused on single-turn refusal---do not measure.

More broadly, this work opens a research direction in \textbf{adversarial agent-to-agent interaction}. As LLM agents are deployed in multi-agent settings---negotiating contracts, managing supply chains, competing in markets---understanding how they behave when their interests conflict becomes critical. Crucible provides a simple, model-agnostic framework for studying these dynamics: swap in any LLM, any game, any prompt condition, and measure what emerges.

We release all code, prompts, game transcripts, metrics, and data at \url{https://github.com/TheApexWu/crucible}.

\bibliography{references}
\bibliographystyle{plainnat}

% ============================================================
\appendix
\section{Agent System Prompts}
\label{app:prompts}
% Full competitive and neutral prompts

\section{Game Transcripts}
\label{app:transcripts}
% Round 1 betrayal transcript, Round 6 reflection, full 10-round game

\section{Private Reflection Examples}
\label{app:reflections}
% Agent A Round 1 (Machiavellian planning), Agent B Round 1 (detection), later-round escalation

\section{Deception Index Derivation}
\label{app:deception-index}
% Weight rationale, sensitivity analysis, calibration

\end{document}
